\documentclass{ctexart}
\usepackage{Note}
\usepackage{unicode-math}
\setCJKmainfont[
  BoldFont = 思源宋体 Bold,
  ItalicFont = 楷体,
]{宋体}
\begin{document}
\section{形式理论}
\subsection{Hilbert空间}
量子理论是建立在两个基础概念上的: \tbf{波函数}和\tbf{算符}.体系的\tbf{状态}由其波函数表示, \tbf{可观察量}用算符来表示.在数学上,我们可以将波函数视作矢量,算符视作作用于矢量上的线性变换,因此量子力学的自然语言是线性代数.\\
\indent 和一般的有限维向量不同,所有波函数的的集合构成了一个无限维的向量空间.当然,为了表示一个可能的物理状态,波函数$\mathbfit{\Psi}$必须是归一化的.为此,考虑在一个特定的区间$[a,b]$(通常是$(-\infty,+\infty)$)内所有平方可积函数的集合.这种空间是一个Hilbert空间.
\begin{definition}[Hilbert空间]
    完备的内积空间称作\tbf{Hilbert空间}.
\end{definition}
前面所述的空间的内积定义如下:
\begin{definition}[波函数空间的内积]
    波函数空间上的内积定义为
    \[\left\langle f|g\right\rangle=\int_{a}^{b}f^\ast(x)g(x)\di x\]
\end{definition}
如果一个函数$f$与自身的内积为$1$,就称该函数是\tbf{归一化的}.如果两个函数的内积为$0$,就称这两个函数是\tbf{正交的}.如果一组函数$\{f_n\}$既是归一的又互相正交,那么称它们\tbf{正交归一}.如果一组函数$\{f_n\}$是Hilbert空间的基(及其延伸组),那么称这组函数是\tbf{完备}的.
\subsection{可观测量}
\subsubsection{Hermitian算符}
引入Hilbert空间及其上的内积后,可观测量$Q(x,p)$的期望值可以简单地用内积符号表示:
\[\avg{Q}=\int\iPsi^\ast\hat{Q}\iPsi\di x=\inprod{\iPsi}{\hat{Q}\iPsi}\]
一次测量的结果应当是实数,因此
\[\avg{Q}=\avg{Q}^\ast\]
而内积的复共轭等于交换参与内积的两个函数形成的内积,因此
\[\inprod{\iPsi}{\hat{Q}\iPsi}=\inprod{\hat{Q}\iPsi}{\iPsi}\]
可观测量对应的算符都应满足上述性质,它们有一个特殊的名字:
\begin{definition}[Hermitian算符]
    对任意波函数$f$都有\[\inprod{f}{\hat{Q}f}=\inprod{\hat{Q}f}{f}\]的算符$\hat{Q}$被称作\tbf{厄米(Hermitian)算符}.
\end{definition}
\subsubsection{定态}
一般而言,对一组全同的系统组成的系综中对可观测量$Q$进行测量时,即使每个系统都处于相同的状态$\iPsi$,每次测量也会得到不同的结果,这就是量子力学的不确定性.\\
\indent 问题在于,是否存在一个状态使得每一次对$Q$的测量都能得到同样的值?换言之,这样的状态$\iPsi$被称作可观测量$Q$的一个\tbf{定态}.在定态下, $Q$的标准差应当是$0$,也就是
\[\sigma^2=\avg{\left(Q-\avg{Q}\right)^2}=\inprod{\iPsi}{\left(\hat{Q}-q\right)^2\iPsi}=\inprod{\left(\hat{Q}-q\right)\iPsi}{\left(\hat{Q}-q\right)\iPsi}=0\]
于是
\[\hat{Q}\iPsi=q\iPsi\]
\begin{definition}[本征值]
    算符$\hat{Q}$的\tbf{本征函数}$\iPsi$是满足
    \[\hat{Q}\iPsi=q\iPsi\]
    的函数, $q$成为对应于该本征函数的\tbf{本征值}.
\end{definition}
\begin{definition}[谱和简并]
    算符的所有本征值的集合称作算符的\tbf{谱}.如果同一本征值对应了线性无关的本征函数,这称作谱的\tbf{简并}.
\end{definition}
\subsection{厄米算符的本征函数}
现在,我们把注意力放在厄米算符的本征函数上,即可观测量具有确定值的状态.这可以分为两种情况:如果谱是离散的,那么本征函数就在Hilbert空间中并且构成物理上可以实现的态;如果谱是连续的,那么本征函数就不可归一化,并且它们无法代表可能实现的波函数.
\subsubsection{离散谱}
在数学上,厄米算符的可归一化的本征函数具有两个重要特征:
\begin{theorem}[厄米算符的本征函数的性质]
    厄米算符的可归一化的本征函数对应的本征值是实数,并且对应于不同本征值的本征函数是正交的.
\end{theorem}
\begin{proof}
    设厄米算符$\hat{Q}$和Hilbert空间内的函数$f$使得
    \[\hat{Q}f=qf\]
    于是
    \[\inprod{f}{\hat{Q}f}=\inprod{\hat{Q}f}{f}\]
    于是
    \[q\inprod ff=q^\ast\inprod ff\]
    因为$f\neq0$,于是$q=q^\ast$,则$q\in\R$.\\
    后一点是容易证明的,可以类比对应于不同特征值的特征向量必然正交来证明.
\end{proof}
在有限维向量空间中,厄米矩阵的特征向量组构成这一空间的基.然而,在无限维的Hilbert空间中,这一性质不能推广,然而这个性质对量子力学内在的自洽性是非常必要的,所以遵从狄拉克的论述,我们给出如下的公理:
\begin{theorem}
    可观测量算符的本征函数是完备的,在Hilbert空间中的任何函数都可以用它们的线性组合来表示.
\end{theorem}
\subsubsection{连续谱}
如果厄米算符的谱是连续的,其本征函数就是不可归一化的,前面的证明也就不能成立.尽管如此,在某种意义上前面三个基本的性质仍然能成立.
\begin{problem}
    在$-\infty<x<\infty$的区间上求动量算符的本征值和本征函数.
\end{problem}
\begin{solution}
    设$f_p(x)$是$\hat{p}$的本征函数,本征值为$p$,则有
    \[-\i\hbar\dfrac{\di}{\di x}f_p(x)=pf_p(x)\]
    这一方程的通解为
    \[f_p(x)=A\exp\left(\dfrac{\i px}{\hbar}\right)\]
    如果$p$是复数,那么上述函数中就会出现$\e^{kx}$项,这会使得$f_p(x)$不能平方可积.看来动量算符在希尔伯特空间中没有本征函数.\\
    然而,如果限定本征值为实数的情况下,就可以有
    \[\int_{-\infty}^{\infty}f^\ast_{p'}(x)f_p(x)\di x=|A|^2\int_{-\infty}^{\infty}\exp\left(\dfrac{\i(p-p')x}{\hbar}\right)\di x=|A|^22\pi\hbar\delta(p-p')\]
    取$A=\dfrac{1}{\sqrt{2\pi\hbar}}$则有
    \[f_p(x)=\dfrac{1}{\sqrt{2\pi\hbar}}{\exp\left(\dfrac{\i px}{\hbar}\right)}\]
    并且
    \[\inprod{f_{p'}}{f_p}=\delta(p-p')\]
    这几乎就是正交归一性的表述,除了把克罗内克$\delta$符号替换成了狄拉克$\delta$函数.这一性质被称作\tbf{狄拉克正交归一性}.\\
    更重要的是,此时动量算符的对应于实数本征值的本征函数是完备的.由于谱的连续性,应当把线性组合替换为积分,这样对于任意平方可积的函数$f(x)$都有
    \[f(x)=\int_{-\infty}^{\infty}c(p)f_p(x)\di p=\dfrac{1}{\sqrt{2\pi\hbar}}\int_{-\infty}^{\infty}c(p)\exp\left(\dfrac{\i px}{\hbar}\right)\di p\]
    并且展开的系数$c(p)$可以通过傅里叶变换得到:
    \[c(p')=\int_{-\infty}^{\infty}c(p)\delta(p-p')\di p=\int_{-\infty}^{\infty}c(p)\inprod{f_{p'}}{f_p}\di p=\inprod{f_{p'}}{f}\]
\end{solution}
这样,我们就得到动量本征函数$f_p$,它是波长$\lambda=\dfrac{2\pi\hbar}{p}$的正弦波.
\subsection{广义统计诠释}
我们在前面的章节介绍了波函数的实际含义,以及可观测量的确定方法.这些内容可以描述为:
\begin{theorem}[广义统计诠释]
    如果对处于$\iPsi(x,t)$状态粒子的可观测量$Q(x,p)$进行测量,那么一定会得到厄米算符$\hat{Q}$的某个本征值.如果$\hat{Q}$的谱是离散的,那么得到与本征函数$\psi_n(x)$相应的本征值$q_n$的概率为$\left|\inprod{\psi_n}{\iPsi}\right|^2$;如果$\hat{Q}$的谱是连续的,并且具有实的本征值$q(z)$,那么结果的概率密度函数为$\rho(z)=\left|\inprod{\psi_z}{\iPsi}\right|^2$.
\end{theorem}
在进行测量后,波函数\tbf{坍缩}为相应的本征态.在连续谱的情况下,坍缩是朝向测量值的一个狭窄的范围,取决于仪器的精度.\\
\indent 总的几率当然也是归一化的.这源于波函数的归一化.对于离散谱有
\[\begin{aligned}
    1
    &=\inprod{\iPsi}{\iPsi}=\inprod{\left(\sum_{n'}c_{n'}f_{n'}\right)}{\left(\sum_{n}c_{n}f_{n}\right)}=\sum_{n'}\sum_{n}c_{n'}^\ast c_{n}\inprod{f_{n'}}{f_n}\\
    &=\sum_{n'}\sum_{n}c_{n'}^\ast c_{n}\delta_{n'n}=\sum_{n}c_n^\ast c_n=\sum_{n}\left|c_n\right|^2
\end{aligned}\]
对于$Q$的期望值也可以得出符合前面章节的结果.\\
\indent 对于连续谱,也可以用广义统计诠释得到我们最初的论述.对于位置$x$,我们知道每个实数$y$都是位置算符$\hat{x}$的本征值,相应的本征函数$g_y(x)=\delta(x-y)$.于是系数$c(y)$为
\[c(y)=\inprod{g_y}{\iPsi}=\int_{-\infty}^{+\infty}\delta(x-y)\iPsi(x,t)\di x=\Psi(y,t)\]
从而概率密度函数$\rho(y)=\left|\iPsi(y,t)\right|^2$.\\
\indent 而对于动量$p$,动量算符$\hat{p}$的本征函数为
\[f_p(x)=\dfrac{1}{\sqrt{2\pi\hbar}}\exp\left(\dfrac{\i px}{\hbar}\right)\]
其系数为
\[c(p)=\inprod{f_p}{\iPsi}=\dfrac{1}{\sqrt{2\pi\hbar}}\int_{-\infty}^{+\infty}\exp\left(-\dfrac{\i px}{\hbar}\right)\iPsi(x,t)\di x\]
这是一个非常重要的物理量:
\begin{definition}[动量空间波函数]
    动量空间波函数$\iPhi(p,t)$为坐标空间波函数$\iPsi(x,t)$的傅里叶变换:
    \[\iPhi(p,t)=\dfrac{1}{\sqrt{2\pi\hbar}}\int_{-\infty}^{+\infty}\exp\left(-\dfrac{\i px}{\hbar}\right)\iPsi(x,t)\di x\]
\end{definition}
同样按照广义统计诠释,对动量进行测量的概率密度函数即为$\left|\iPhi(p,t)\right|^2$.
\subsection{不确定原理}
在严格叙述广义统计诠释后,我们现在来推导(广义的)不确定原理.
\begin{theorem}[不确定原理]
    对于可观测量$A$和$B$,总有
    \[\sigma_A^2\sigma_B^2\geq\left(\dfrac{1}{2\i}\avg{\left[\hat{A},\hat{B}\right]}\right)^2\]
\end{theorem}
\begin{proof}
    根据方差的定义有
    \[\sigma_A^2=\avg{(\avg{A}-A)^2}=\inprod{\left(\hat{A}-\avg{A}\right)\iPsi}{\left(\hat{A}-\avg{A}\right)\iPsi}=\inprod{f}{f},\quad f=\left(\hat{A}-\avg{A}\right)\iPsi\]
    同样地有
    \[\sigma_B^2=\inprod gg,\quad g=\left(\hat{B}-\avg{B}\right)\iPsi\]
    根据Cauchy-Schwarz不等式有
    \[\sigma_A^2\sigma_B^2=\inprod ff\inprod gg\geq\left|\inprod fg\right|^2\]
    由于对于任意复数$z$有
    \[\left|z^2\right|=\left[\Re(z)\right]^2+\left[\Im(z)\right]^2\geq\left[\Im(z)\right]^2=\left[\dfrac{1}{2\i}\left(z-z^\ast\right)\right]^2\]
    从而令$z=\inprod fg$就有
    \[\sigma_A^2\sigma_B^2\geq\left(\dfrac{1}{2\i}\left[\inprod fg-\inprod gf\right]\right)^2\]
    根据厄米算符的性质有
    \[\begin{aligned}
        \inprod fg
        &=\inprod{\left(\hat{A}-\avg{A}\right)\iPsi}{\left(\hat{B}-\avg{B}\right)\iPsi}=\inprod{\iPsi}{\left(\hat{A}-\avg{A}\right)\left(\hat{B}-\avg{B}\right)\iPsi}\\
        &=\inprod{\iPsi}{\hat{A}\hat{B}\iPsi}-\avg{A}\inprod{\iPsi}{\hat{B}\iPsi}-\avg{B}\inprod{\iPsi}{\hat{A}\iPsi}+\avg{A}\avg{B}\inprod{\iPsi}{\iPsi}\\
        &=\avg{\hat{A}\hat{B}}-\avg{A}\avg{B}
    \end{aligned}\]
    同理有
    \[\inprod gf=\avg{\hat{B}\hat{A}}-\avg{A}\avg{B}\]
    于是
    \[\inprod fg-\inprod gf=\avg{\hat{A}\hat{B}}-\avg{\hat{B}\hat{A}}=\avg{\left[\hat{A},\hat{B}\right]}\]
    于是
    \[\sigma_A^2\sigma_B^2\geq\left(\dfrac{1}{2\i}\avg{\left[\hat{A},\hat{B}\right]}\right)^2\]
\end{proof}
事实上,对于任意一对可观测量,如果其算符不对易,那么这两个观测量就存在不确定原理,我们称这对可观测量为\tbf{不相容可观测量}.不相容可观测量没有相同的本征函数(或者更严谨地说,不存在相同的本征函数构成的完备集).\\
\indent 除去比较著名的不确定原理,即位置-动量不确定原理外,另一对常见的不相容的可观测量是时间$t$和能量$E$.
\begin{theorem}[能量-时间不确定原理]
    能量$E$和时间$t$测定的标准差满足
    \[\Delta t\Delta E\geq\dfrac{\hbar}{2}\]
\end{theorem}
在相对论中,这一结论是比较容易说明的.但是我们现在希望从非相对论的量子力学理论中推导出此式;并且事实上,尽管它与位置-动量不确定原理的形式相似,原理却不同.毕竟,位置$x$,动量$p$和能量$E$都是系统的动力学变量,但时间$t$却不是.
\begin{proof}
    时间$t$是衡量体系变化快慢的量度.我们考虑可观测量$Q(x,p,t)$的期望对时间$t$的导数:
    \[\dfrac{\d}{\d t}\avg{Q}=\dfrac{\d}{\d t}\inprod{\iPsi}{\hat{Q}\iPsi}=\inprod{\dfrac{\p\iPsi}{\p t}}{\hat{Q}\iPsi}+\inprod{\iPsi}{\dfrac{\p\hat{Q}}{\p t}\iPsi}+\inprod{\iPsi}{\hat{Q}\dfrac{\p\iPsi}{\p t}}\]
    根据薛定谔方程
    \[\i\hbar\dfrac{\p\iPsi}{\p t}=\hat{H}\iPsi\]
    有
    \[\dfrac{\d}{\d t}\avg{Q}=-\dfrac{1}{\i\hbar}\inprod{\hat{H}\iPsi}{\hat{Q}\iPsi}+\dfrac{1}{\i\hbar}\inprod{\iPsi}{\hat{Q}\hat{H}\iPsi}+\avg{\dfrac{\p\hat{Q}}{\p t}}\]
    而$\hat{H}$本身是厄米算符,于是
    \[\dfrac{\d}{\d t}\avg{Q}=\dfrac{\i}{\hbar}\avg{\left[\hat{H},\hat{Q}\right]}+\avg{\dfrac{\p\hat{Q}}{\p t}}\]
    这一式子本身是一个重要的公式:
    \begin{theorem}[广义埃伦菲斯特定理]
        可观测量$Q$的期望随时间的变化率为
        \[\dfrac{\d}{\d t}\avg{Q}=\dfrac{\i}{\hbar}\avg{\left[\hat{H},\hat{Q}\right]}+\avg{\dfrac{\p\hat{Q}}{\p t}}\]
    \end{theorem}
    可以看出,在算符不含时的情况下(一般而言都是如此),其期望值的变化率与该算符和哈密顿算符的对易决定.特别地,当$\hat{Q}$与$\hat{H}$对易时,$\avg{Q}$是常量,即一个守恒量.\\
    \indent 现在,我们在广义不确定关系中令$A=H$和$B=Q$,并且假定$Q$不含时,则有
    \[\sigma_H^2\sigma_Q^2\geq\left(\dfrac{1}{2\i}\avg{\left[\hat{H},\hat{Q}\right]}\right)^2=\left(\dfrac{1}{2\i}\dfrac{\hbar}{\i}\dfrac{\d\avg{Q}}{\d t}\right)^2\]
    即
    \[\sigma_H\sigma_Q\geq\dfrac{\hbar}{2}\left|\dfrac{\d\avg{Q}}{\d t}\right|\]
    定义$\Delta E=\sigma_H$以及$\Delta t=\dfrac{\sigma_Q}{\left|\d\avg{Q}/\d t\right|}$,则有
    \[\Delta E\Delta t\geq\dfrac{\hbar}{2}\]
    这里$\Delta t$表示$Q$的期望值大小变化一个标准差所需的时间.特别地, $\Delta t$总是依赖于可观测量$Q$.不过,当$\Delta E$比较小时,所有的可观测量的变化都是比较平缓的;反之,如果某个可观测量变化很快,那么能量的不确定性就非常大.
\end{proof}
\subsection{矢量和算符}
\subsubsection{Hilbert空间的基矢}
一个向量空间中的矢量$\vec{u}$可以用不同的基进行表达,尽管表达的方式不同,但它们都指向同一个矢量.量子力学中体系的状态也是如此,它由Hilbert空间中的一个矢量来描述: $\ket{\mathcal{S}(t)}$\footnote{这是被称作Dirac符号中的\tbf{右矢}.我们将马上正式地介绍Dirac符号.}.那么,波函数$\iPsi(x,t)$实际上是状态$\ket{\mathcal{S}(t)}$在坐标本征函数作为基上展开的分量:
\[\iPsi(x,t)=\inprod{x}{\mathcal{S}(t)}\]
这里基矢量$\ket{x}$是本征值为$x$的算符$\hat{x}$的本征函数,求内积时取其共轭转置$\bra{x}$.同样地,动量空间上的波函数是$\ket{\mathcal{S}(t)}$以动量本征函数为基进行展开时的$p$分量:
\[\iPhi(p,t)=\inprod{p}{\mathcal{S}(t)}\]
$\bra{p}$的含义同理.或者我们可以把$\ket{\mathcal{S}(t)}$用有离散谱的算符(例如能量)的本征函数的基展开:
\[c_n(t)=\inprod{n}{\mathcal{S}(t)}\]
这里基矢量$\bra{n}$是$\hat{H}$的第$n$个本征函数.所有上面的表示方法都是同一个状态,表示一样的信息,仅仅是描述同一状态(矢量)的在三组基下的展开方式:
\[\ket{\mathcal{S}(t)}\to\int\iPsi(y,t)\delta(x-y)\di y=\int\iPhi(p,t)\dfrac{1}{\sqrt{2\pi\hbar}}\exp\left(\dfrac{\i px}{\hbar}\right)\d p=\sum_{n}c_n\exp\left(-\dfrac{\i E_n t}{\hbar}\right)\psi_n(x)\]
\indent 而算符是Hilbert空间上的线性变换:
\[\ket{\beta}=\hat{Q}\ket\alpha\]
算符$\hat{Q}$所对应的矩阵元表示了分量的变换情况,这在线性代数中已经学习过.\\
\indent 尽管状态$\ket{\mathcal{S}(t)}$是Hilbert空间上的矢量,但实际情况下的大多数问题都是在坐标空间中讨论的,因此把波函数$\iPsi(x,t)$称作\textit{体系的状态}并无大碍.
\subsubsection{Dirac符号}
Dirac提议把内积分成两部分:
\begin{definition}[Dirac符号]
    内积$\inprod{\alpha}{\beta}$可以分为两部分:\tbf{左矢(bra)} $\bra{\alpha}$和\tbf{右矢(ket)} $\ket(\beta)$.
\end{definition}
右矢代表一个矢量,而左矢作为右矢的共轭转置实际上也是一种算符,其组成的空间是原来空间的\tbf{对偶空间}.这样的记号提供了一种简洁的数学工具.例如:
\begin{definition}[投影算符]
    归一化矢量$\ket{\alpha}$对应的\tbf{投影算符}为:
    \[\hat{P}=\ket{\alpha}\bra{\alpha}\]
    对于任意的向量$\ket{\beta}$,投影算符给出其在$\ket{\alpha}$上的投影向量:
    \[\hat{P}\ket{\beta}=\inprod{\alpha}{\beta}\ket{\alpha}\]
\end{definition}
并且如果$\left\{\ket{e_n}\right\}$是离散的规范正交基,则有
\[\sum_{n}\ket{e_n}\bra{e_n}=\mbf{1}\]
为\tbf{恒等算符}.这样,对于任意的向量$\ket{\alpha}$都有
\[\alpha=\mbf1\alpha=\sum_{n}\ket{e_n}\inprod{e_n}{\alpha}\]
类似地,如果$\left\{\ket{e_z}\right\}$是连续的Dirac正交归一的连续基,那么
\[\int\ket{e_z}\bra{e_z}\di z=\mbf{1}\]
\subsubsection{Dirac符号中的基矢变换}
Dirac符号的优点是我们讨论问题可以不局限在某一组特定的基上,而且可以方便地在不同的基之间进行变换.这里以一个引理的证明为例.
\begin{lemma}
    动量算符$\hat{p}$的形式可以由正则对易关系$\left[\hat{x},\hat{p}\right]=\i\hbar$推出.
\end{lemma}
\begin{proof}
    
\end{proof}
\end{document}